\documentclass[11pt,a4paper]{article}
\usepackage[utf8]{inputenc}
\usepackage[T1]{fontenc}
\usepackage[catalan]{babel}
\usepackage{lmodern}
\usepackage[sfdefault,lining]{FiraSans}
\usepackage[a4paper,top=2.3cm,bottom=2.2cm,left=1.8cm,right=1.8cm,headheight=1.3cm,headsep=0.4cm,footskip=1cm]{geometry}
\usepackage{amsmath,amssymb}
\usepackage{icomma}
\usepackage{graphicx}
\usepackage[table]{xcolor}
\usepackage{tikz}
\usepackage{fontawesome5}
\usepackage{ccicons}
\usepackage{fancyhdr}
\usepackage{lastpage}
\usepackage{tasks}
\usepackage{needspace}
\usepackage[most]{tcolorbox}
\usepackage{array}
\usepackage[hidelinks,pdftitle={Solucions de la fitxa de la unitat 4: Equacions. 4t d'ESO Matemàtiques A},pdfsubject={Matemàtiques A 4t ESO},pdfkeywords={equacions, equacions de primer grau, denominadors, problemes, equacions de segon grau, completes, incompletes, discriminant, equacions factoritzades, equacions racionals, equacions amb radicals}]{hyperref}

\definecolor{navy}{HTML}{1B2748}
\definecolor{accent}{HTML}{E8603A}
\definecolor{teal}{HTML}{1F9E90}
\definecolor{soft}{HTML}{EEF1F8}
\definecolor{muted}{HTML}{5D6685}

\setlength{\parindent}{0pt}
\setlength{\parskip}{3pt}
\settasks{label=\alph*), label-format=\bfseries\color{navy}, label-width=1.3em, item-indent=1.8em, column-sep=0.6em, after-item-skip=5pt, before-skip=3pt, after-skip=4pt}

% ---------- capçalera i peu ----------
\newcommand{\llicencia}{\href{https://creativecommons.org/licenses/by-nc-sa/4.0/deed.ca}{\ccbyncsa}\ \ Llicència Creative Commons Reconeixement-NoComercial-CompartirIgual 4.0 Internacional (CC BY-NC-SA 4.0)}
\pagestyle{fancy}
\fancyhf{}
\fancyhead[L]{\raisebox{-0.35\height}{\includegraphics[height=0.95cm]{logo-gva.png}}}
\fancyhead[C]{\small\color{muted}Solucions de la fitxa · Unitat 4: Equacions · 4t ESO A}
\fancyhead[R]{\raisebox{-0.35\height}{\includegraphics[height=0.95cm]{logo-vermellar.png}}}
\fancyfoot[L]{\scriptsize\color{muted}\llicencia}
\fancyfoot[R]{\small\color{muted}\thepage}
\renewcommand{\headrulewidth}{0.4pt}
\renewcommand{\headrule}{\color{muted!40}\hrule height \headrulewidth}
\fancypagestyle{primera}{\fancyhf{}\renewcommand{\headrulewidth}{0pt}
  \fancyfoot[L]{\scriptsize\color{muted}\llicencia}
  \fancyfoot[R]{\small\color{muted}\thepage}}

% ---------- nivells i exercicis ----------
\newcommand{\estrelles}[1]{{\color{accent}\ifcase#1\or\faStar\faStar[regular]\faStar[regular]\or\faStar\faStar\faStar[regular]\or\faStar\faStar\faStar\fi}}
\newcounter{exer}
\newcommand{\ex}[1]{\par\needspace{4\baselineskip}\medskip\refstepcounter{exer}\noindent{\color{navy}\bfseries\theexer.}\ \ {\small\estrelles{#1}}\ \ }

\newcommand{\seccio}[2]{\par\needspace{16\baselineskip}\bigskip
  \begin{tcolorbox}[colback=navy,colframe=navy,coltext=white,arc=2mm,boxrule=0pt,left=3mm,right=3mm,top=1.2mm,bottom=1.2mm]
  {\large\bfseries #1\quad #2}
  \end{tcolorbox}\vspace{-2pt}}

\newcommand{\exf}[1]{\par\medskip\refstepcounter{exer}\noindent{\color{navy}\bfseries\theexer.}\ \ {\small\estrelles{#1}}\ \ }
\newcolumntype{C}[1]{>{\centering\arraybackslash}p{#1}}
% quadre de recordatori (només a la fitxa d'exercicis)
\newcommand{\recorda}[1]{\par\smallskip\begin{tcolorbox}[colback=accent!7,colframe=accent!55,arc=2mm,boxrule=0.5pt,left=2.5mm,right=2.5mm,top=1mm,bottom=1mm]\small{\color{accent}\faLightbulb[regular]}\ \textbf{Recorda:}\ #1\end{tcolorbox}}
\newcommand{\res}[1]{\boxed{#1}}\newcommand{\rest}[1]{\fbox{#1}}

\def\fmneed{11}% que el títol d'una secció no quede sol al final del full
% =====================================================================
% Capçalera FentMates de les fitxes: un color per a cada curs.
% Els colors són els de la pàgina inicial (LEVELS a assets/js/materials.js);
% el logo ve de _eines/logo/pdf/ (python3 _eines/logo/genera.py).
%
% Ús, al final del preàmbul de la unitat:
%   % =====================================================================
% Capçalera FentMates de les fitxes: un color per a cada curs.
% Els colors són els de la pàgina inicial (LEVELS a assets/js/materials.js);
% el logo ve de _eines/logo/pdf/ (python3 _eines/logo/genera.py).
%
% Ús, al final del preàmbul de la unitat:
%   % =====================================================================
% Capçalera FentMates de les fitxes: un color per a cada curs.
% Els colors són els de la pàgina inicial (LEVELS a assets/js/materials.js);
% el logo ve de _eines/logo/pdf/ (python3 _eines/logo/genera.py).
%
% Ús, al final del preàmbul de la unitat:
%   \input{../comu/fentmates.tex}
%   \fmcurs{1eso}
% i, al començament del document:
%   \capcalera{Matemàtiques}{1r ESO}{Fitxa d'exercicis · Unitat 2}{Divisibilitat}
% Opcional abans de \input: \def\fmneed{14} (espai mínim abans de cada secció).
% =====================================================================
\makeatletter
\def\fm@logodir{../../logo/pdf}
\providecommand{\fmneed}{10}
\def\fm@cursnom{}

% clau del curs, color i color del text damunt del color
\newcommand{\fm@curs}[3]{\@namedef{fm@color@#1}{#2}\@namedef{fm@text@#1}{#3}}
\fm@curs{1eso}{E6B84A}{navy}
\fm@curs{2eso}{3B82F6}{white}
\fm@curs{3eso}{FF7A50}{white}
\fm@curs{4eso-a}{A78BFA}{white}
\fm@curs{4eso-b}{6D28D9}{white}
\fm@curs{1batx-mi}{9B1C1C}{white}
\fm@curs{1batx-ccss}{EF6B6B}{white}
\fm@curs{2batx-mii}{1E6B47}{white}
\fm@curs{2batx-ccss}{2BB3A3}{white}

\newcommand{\fmcurs}[1]{%
  \@ifundefined{fm@color@#1}{\PackageError{fentmates}{Curs desconegut: #1}{}}{}%
  \edef\fm@tmp{\noexpand\definecolor{curs}{HTML}{\@nameuse{fm@color@#1}}%
               \noexpand\colorlet{curstext}{\@nameuse{fm@text@#1}}}\fm@tmp
  \colorlet{cursfosc}{curs!72!black}%
  \def\fm@logo{\fm@logodir/vertical-#1.pdf}%
  \def\fm@marca{\fm@logodir/marca-#1.pdf}}

% ---------- capçalera de la primera pàgina ----------
% #1 matèria (petit), #2 curs, #3 tipus de document i unitat, #4 títol
\newcommand{\capcalera}[4]{%
  \gdef\fm@cursnom{#2}%
  \par\noindent
  \begin{tcolorbox}[colback=curs,colframe=curs,boxrule=0pt,arc=2.2mm,width=3.5cm,height=2.15cm,
      valign=center,left=3mm,right=2mm,top=1mm,bottom=1mm,nobeforeafter,box align=center]
    \color{curstext}{\footnotesize #1}\par\vspace{1pt}{\fontsize{19}{21}\selectfont\bfseries #2\par}%
  \end{tcolorbox}\hspace{5mm}%
  \begin{minipage}[c]{\dimexpr\linewidth-3.5cm-5mm-2.9cm\relax}
    {\color{cursfosc}\small\bfseries\MakeUppercase{#3}}\par\vspace{2pt}
    {\color{navy}\fontsize{20}{23}\selectfont\bfseries #4\par}
  \end{minipage}\hfill
  \begin{minipage}[c]{2.6cm}\raggedleft\includegraphics[width=2.6cm]{\fm@logo}\end{minipage}\par
  \vspace{2pt}}

% ---------- barres de secció del color del curs ----------
\renewcommand{\seccio}[2]{\par\needspace{\fmneed\baselineskip}\bigskip
  \begin{tcolorbox}[colback=curs,colframe=curs,coltext=curstext,arc=2mm,boxrule=0pt,left=3mm,right=3mm,top=1.2mm,bottom=1.2mm]
  {\large\bfseries #1\quad #2}
  \end{tcolorbox}\vspace{-2pt}}

% ---------- estrelles de dificultat del color del curs (a Batxillerat no se'n posen) ----------
\@ifundefined{estrelles}{}{%
  \renewcommand{\estrelles}[1]{{\color{curs}\ifcase#1\or\faStar\faStar[regular]\faStar[regular]\or\faStar\faStar\faStar[regular]\or\faStar\faStar\faStar\fi}}}

% ---------- capçalera i peu de les pàgines ----------
% primera pàgina: logos institucionals al peu
\fancypagestyle{primera}{\fancyhf{}\renewcommand{\headrulewidth}{0pt}
  \fancyfoot[L]{\raisebox{-0.32\height}{\includegraphics[height=0.78cm]{logo-gva.png}}\hspace{5mm}%
    \raisebox{-0.32\height}{\includegraphics[height=0.78cm]{logo-vermellar.png}}}
  \fancyfoot[R]{\scriptsize\color{muted}\href{https://creativecommons.org/licenses/by-nc-sa/4.0/deed.ca}{\ccbyncsa}\ \ CC BY-NC-SA 4.0 · fentmates.org\qquad\small\thepage}}
% la resta: el símbol a l'esquerra i el curs a la dreta, amb el títol al centre com abans
\fancyhead[L]{\raisebox{-0.3\height}{\includegraphics[height=0.62cm]{\fm@marca}}}
\fancyhead[R]{\small\bfseries\color{cursfosc}\fm@cursnom}
\renewcommand{\headrule}{\color{curs!45}\hrule height \headrulewidth}
\makeatother

%   \fmcurs{1eso}
% i, al començament del document:
%   \capcalera{Matemàtiques}{1r ESO}{Fitxa d'exercicis · Unitat 2}{Divisibilitat}
% Opcional abans de \input: \def\fmneed{14} (espai mínim abans de cada secció).
% =====================================================================
\makeatletter
\def\fm@logodir{../../logo/pdf}
\providecommand{\fmneed}{10}
\def\fm@cursnom{}

% clau del curs, color i color del text damunt del color
\newcommand{\fm@curs}[3]{\@namedef{fm@color@#1}{#2}\@namedef{fm@text@#1}{#3}}
\fm@curs{1eso}{E6B84A}{navy}
\fm@curs{2eso}{3B82F6}{white}
\fm@curs{3eso}{FF7A50}{white}
\fm@curs{4eso-a}{A78BFA}{white}
\fm@curs{4eso-b}{6D28D9}{white}
\fm@curs{1batx-mi}{9B1C1C}{white}
\fm@curs{1batx-ccss}{EF6B6B}{white}
\fm@curs{2batx-mii}{1E6B47}{white}
\fm@curs{2batx-ccss}{2BB3A3}{white}

\newcommand{\fmcurs}[1]{%
  \@ifundefined{fm@color@#1}{\PackageError{fentmates}{Curs desconegut: #1}{}}{}%
  \edef\fm@tmp{\noexpand\definecolor{curs}{HTML}{\@nameuse{fm@color@#1}}%
               \noexpand\colorlet{curstext}{\@nameuse{fm@text@#1}}}\fm@tmp
  \colorlet{cursfosc}{curs!72!black}%
  \def\fm@logo{\fm@logodir/vertical-#1.pdf}%
  \def\fm@marca{\fm@logodir/marca-#1.pdf}}

% ---------- capçalera de la primera pàgina ----------
% #1 matèria (petit), #2 curs, #3 tipus de document i unitat, #4 títol
\newcommand{\capcalera}[4]{%
  \gdef\fm@cursnom{#2}%
  \par\noindent
  \begin{tcolorbox}[colback=curs,colframe=curs,boxrule=0pt,arc=2.2mm,width=3.5cm,height=2.15cm,
      valign=center,left=3mm,right=2mm,top=1mm,bottom=1mm,nobeforeafter,box align=center]
    \color{curstext}{\footnotesize #1}\par\vspace{1pt}{\fontsize{19}{21}\selectfont\bfseries #2\par}%
  \end{tcolorbox}\hspace{5mm}%
  \begin{minipage}[c]{\dimexpr\linewidth-3.5cm-5mm-2.9cm\relax}
    {\color{cursfosc}\small\bfseries\MakeUppercase{#3}}\par\vspace{2pt}
    {\color{navy}\fontsize{20}{23}\selectfont\bfseries #4\par}
  \end{minipage}\hfill
  \begin{minipage}[c]{2.6cm}\raggedleft\includegraphics[width=2.6cm]{\fm@logo}\end{minipage}\par
  \vspace{2pt}}

% ---------- barres de secció del color del curs ----------
\renewcommand{\seccio}[2]{\par\needspace{\fmneed\baselineskip}\bigskip
  \begin{tcolorbox}[colback=curs,colframe=curs,coltext=curstext,arc=2mm,boxrule=0pt,left=3mm,right=3mm,top=1.2mm,bottom=1.2mm]
  {\large\bfseries #1\quad #2}
  \end{tcolorbox}\vspace{-2pt}}

% ---------- estrelles de dificultat del color del curs (a Batxillerat no se'n posen) ----------
\@ifundefined{estrelles}{}{%
  \renewcommand{\estrelles}[1]{{\color{curs}\ifcase#1\or\faStar\faStar[regular]\faStar[regular]\or\faStar\faStar\faStar[regular]\or\faStar\faStar\faStar\fi}}}

% ---------- capçalera i peu de les pàgines ----------
% primera pàgina: logos institucionals al peu
\fancypagestyle{primera}{\fancyhf{}\renewcommand{\headrulewidth}{0pt}
  \fancyfoot[L]{\raisebox{-0.32\height}{\includegraphics[height=0.78cm]{logo-gva.png}}\hspace{5mm}%
    \raisebox{-0.32\height}{\includegraphics[height=0.78cm]{logo-vermellar.png}}}
  \fancyfoot[R]{\scriptsize\color{muted}\href{https://creativecommons.org/licenses/by-nc-sa/4.0/deed.ca}{\ccbyncsa}\ \ CC BY-NC-SA 4.0 · fentmates.org\qquad\small\thepage}}
% la resta: el símbol a l'esquerra i el curs a la dreta, amb el títol al centre com abans
\fancyhead[L]{\raisebox{-0.3\height}{\includegraphics[height=0.62cm]{\fm@marca}}}
\fancyhead[R]{\small\bfseries\color{cursfosc}\fm@cursnom}
\renewcommand{\headrule}{\color{curs!45}\hrule height \headrulewidth}
\makeatother

%   \fmcurs{1eso}
% i, al començament del document:
%   \capcalera{Matemàtiques}{1r ESO}{Fitxa d'exercicis · Unitat 2}{Divisibilitat}
% Opcional abans de \input: \def\fmneed{14} (espai mínim abans de cada secció).
% =====================================================================
\makeatletter
\def\fm@logodir{../../logo/pdf}
\providecommand{\fmneed}{10}
\def\fm@cursnom{}

% clau del curs, color i color del text damunt del color
\newcommand{\fm@curs}[3]{\@namedef{fm@color@#1}{#2}\@namedef{fm@text@#1}{#3}}
\fm@curs{1eso}{E6B84A}{navy}
\fm@curs{2eso}{3B82F6}{white}
\fm@curs{3eso}{FF7A50}{white}
\fm@curs{4eso-a}{A78BFA}{white}
\fm@curs{4eso-b}{6D28D9}{white}
\fm@curs{1batx-mi}{9B1C1C}{white}
\fm@curs{1batx-ccss}{EF6B6B}{white}
\fm@curs{2batx-mii}{1E6B47}{white}
\fm@curs{2batx-ccss}{2BB3A3}{white}

\newcommand{\fmcurs}[1]{%
  \@ifundefined{fm@color@#1}{\PackageError{fentmates}{Curs desconegut: #1}{}}{}%
  \edef\fm@tmp{\noexpand\definecolor{curs}{HTML}{\@nameuse{fm@color@#1}}%
               \noexpand\colorlet{curstext}{\@nameuse{fm@text@#1}}}\fm@tmp
  \colorlet{cursfosc}{curs!72!black}%
  \def\fm@logo{\fm@logodir/vertical-#1.pdf}%
  \def\fm@marca{\fm@logodir/marca-#1.pdf}}

% ---------- capçalera de la primera pàgina ----------
% #1 matèria (petit), #2 curs, #3 tipus de document i unitat, #4 títol
\newcommand{\capcalera}[4]{%
  \gdef\fm@cursnom{#2}%
  \par\noindent
  \begin{tcolorbox}[colback=curs,colframe=curs,boxrule=0pt,arc=2.2mm,width=3.5cm,height=2.15cm,
      valign=center,left=3mm,right=2mm,top=1mm,bottom=1mm,nobeforeafter,box align=center]
    \color{curstext}{\footnotesize #1}\par\vspace{1pt}{\fontsize{19}{21}\selectfont\bfseries #2\par}%
  \end{tcolorbox}\hspace{5mm}%
  \begin{minipage}[c]{\dimexpr\linewidth-3.5cm-5mm-2.9cm\relax}
    {\color{cursfosc}\small\bfseries\MakeUppercase{#3}}\par\vspace{2pt}
    {\color{navy}\fontsize{20}{23}\selectfont\bfseries #4\par}
  \end{minipage}\hfill
  \begin{minipage}[c]{2.6cm}\raggedleft\includegraphics[width=2.6cm]{\fm@logo}\end{minipage}\par
  \vspace{2pt}}

% ---------- barres de secció del color del curs ----------
\renewcommand{\seccio}[2]{\par\needspace{\fmneed\baselineskip}\bigskip
  \begin{tcolorbox}[colback=curs,colframe=curs,coltext=curstext,arc=2mm,boxrule=0pt,left=3mm,right=3mm,top=1.2mm,bottom=1.2mm]
  {\large\bfseries #1\quad #2}
  \end{tcolorbox}\vspace{-2pt}}

% ---------- estrelles de dificultat del color del curs (a Batxillerat no se'n posen) ----------
\@ifundefined{estrelles}{}{%
  \renewcommand{\estrelles}[1]{{\color{curs}\ifcase#1\or\faStar\faStar[regular]\faStar[regular]\or\faStar\faStar\faStar[regular]\or\faStar\faStar\faStar\fi}}}

% ---------- capçalera i peu de les pàgines ----------
% primera pàgina: logos institucionals al peu
\fancypagestyle{primera}{\fancyhf{}\renewcommand{\headrulewidth}{0pt}
  \fancyfoot[L]{\raisebox{-0.32\height}{\includegraphics[height=0.78cm]{logo-gva.png}}\hspace{5mm}%
    \raisebox{-0.32\height}{\includegraphics[height=0.78cm]{logo-vermellar.png}}}
  \fancyfoot[R]{\scriptsize\color{muted}\href{https://creativecommons.org/licenses/by-nc-sa/4.0/deed.ca}{\ccbyncsa}\ \ CC BY-NC-SA 4.0 · fentmates.org\qquad\small\thepage}}
% la resta: el símbol a l'esquerra i el curs a la dreta, amb el títol al centre com abans
\fancyhead[L]{\raisebox{-0.3\height}{\includegraphics[height=0.62cm]{\fm@marca}}}
\fancyhead[R]{\small\bfseries\color{cursfosc}\fm@cursnom}
\renewcommand{\headrule}{\color{curs!45}\hrule height \headrulewidth}
\makeatother

\fmcurs{4eso-a}

\begin{document}
\thispagestyle{primera}

\capcalera{Matemàtiques A}{4t ESO}{Solucions de la fitxa · Unitat 4}{\raggedright Equacions}
\vspace{8pt}

\begin{tcolorbox}[colback=soft,colframe=soft,arc=2mm,left=3mm,right=3mm,top=1.5mm,bottom=1.5mm]
\small Intenta resoldre cada exercici abans de mirar-ne la solució. Si no t'ix el mateix resultat, repassa els passos o consulta els apunts interactius de la unitat.
\end{tcolorbox}

% =====================================================
\seccio{1}{Equacions de primer grau}

\exf{1} \begin{tasks}(3)
  \task $\res{x=12}$
  \task $3x=-15\Rightarrow$ $\res{x=-5}$
  \task $\res{x=10}$
  \task $-4x=-8\Rightarrow$ $\res{x=2}$
  \task $4x=12\Rightarrow$ $\res{x=3}$
  \task $-3x=-6\Rightarrow$ $\res{x=2}$
\end{tasks}

\ex{2} \begin{tasks}(2)
  \task $-3x=-3\Rightarrow$ $\res{x=1}$
  \task $-3x=10\Rightarrow$ $\res{x=-\frac{10}{3}}$
  \task $-3x=1\Rightarrow$ $\res{x=-\frac{1}{3}}$
  \task $0x=0$: $\res{\text{infinites solucions}}$
  \task $0x=-2$: $\res{\text{no té solució}}$
  \task $-4x=-14\Rightarrow$ $\res{x=\frac{7}{2}}$
\end{tasks}

\ex{2} \begin{tasks}(2)
  \task $\cdot10$: $-4x=-16\Rightarrow$ $\res{x=4}$
  \task $\cdot10$: $10x=60\Rightarrow$ $\res{x=6}$
  \task $\cdot10$: $2x=16\Rightarrow$ $\res{x=8}$
\end{tasks}

\ex{2} \begin{tasks}(2)
  \task $6\neq 8$: $\res{\text{no}}$
  \task $10=10$: $\res{\text{sí}}$
  \task $-\frac{1}{3}\neq -3$: $\res{\text{no}}$
  \task $10=10$: $\res{\text{sí}}$
\end{tasks}

\ex{3} Per exemple: $2(x+5)=x+7$; \ $2(x+1)=2x+5$; \ $3(x-1)=3x-3$.


% =====================================================
\seccio{2}{Equacions de primer grau amb denominadors}

\exf{1} \begin{tasks}(2)
  \task m.c.m.$=6$: $5x=60\Rightarrow$ $\res{x=12}$
  \task m.c.m.$=4$: $3x=36\Rightarrow$ $\res{x=12}$
  \task m.c.m.$=6$: $-x=-5\Rightarrow$ $\res{x=5}$
  \task m.c.m.$=12$: $\res{x=6}$
\end{tasks}

\ex{2} \begin{tasks}(2)
  \task m.c.m.$=12$: $\res{x=23}$
  \task m.c.m.$=10$: $3x=21\Rightarrow$ $\res{x=7}$
  \task m.c.m.$=6$: $\res{x=3}$
  \task m.c.m.$=4$: $2x=2\Rightarrow$ $\res{x=1}$
\end{tasks}

\ex{3} \begin{tasks}(2)
  \task m.c.m.$=12$: $-3x=-11\Rightarrow$ $\res{x=\frac{11}{3}}$
  \task m.c.m.$=6$: $0x=-4$: $\res{\text{no té solució}}$
  \task m.c.m.$=4$: $0x=0$: $\res{\text{infinites solucions}}$
  \task m.c.m.$=12$: $4x=11\Rightarrow$ $\res{x=\frac{11}{4}}$
\end{tasks}

\ex{3} \begin{tasks}(2)
  \task m.c.m.$=20$: $2x=-5\Rightarrow$ $\res{x=-\frac{5}{2}}$
  \task m.c.m.$=12$: $0x=-7$: $\res{\text{no té solució}}$
\end{tasks}


% =====================================================
\seccio{3}{Problemes de primer grau}

\exf{1} $x+(x+1)+(x+2)=147\Rightarrow3x=144\Rightarrow x=48$: $\res{48,\ 49,\ 50}$

\ex{1} $0,8x=68\Rightarrow x=\res{85\ \text{€}}$.

\ex{2} $41-x=4(14-x)\Rightarrow3x=15\Rightarrow x=5$: $\res{\text{fa 5 anys}}$ (36 i 9 anys).

\ex{2} $2x+(x-4)=50\Rightarrow x=18$: $\res{18,\ 18\ \text{i}\ 14\ \text{cm}}$

\ex{2} Després de dilluns queden $\frac23x$; dimarts se'n gasten $\frac25\cdot\frac23x=\frac4{15}x$: $x-\dfrac x3-\dfrac{4x}{15}=240\Rightarrow\dfrac{6x}{15}=240\Rightarrow x=\res{600\ \text{L}}$.

\ex{2} $8\cdot30+12x=9,5(30+x)\Rightarrow240+12x=285+9,5x\Rightarrow2,5x=45\Rightarrow x=\res{18\ \text{kg}}$.

\ex{2} $x+2x+(2x+200)=1400\Rightarrow5x=1200\Rightarrow x=240$: $\res{240\ \text{€},\ 480\ \text{€}\ \text{i}\ 680\ \text{€}}$

\ex{3} Si el segon tarda $x$ hores, $120x=90(x+1,5)\Rightarrow30x=135\Rightarrow x=4,5$: $\res{4\ \text{h}\ 30\ \text{min}}$, a $\res{540\ \text{km}}$.

\ex{3} En una hora omplin $\frac16$ i $\frac13$ del dipòsit. En $x$ hores: $\dfrac x6+\dfrac x3=1\Rightarrow x+2x=6\Rightarrow x=\res{2\ \text{h}}$.


% =====================================================
\seccio{4}{Equacions de segon grau incompletes}

\exf{1} \begin{tasks}(3)
  \task $x^2=49$: $\res{x=\pm7}$
  \task $x^2=9$: $\res{x=\pm3}$
  \task $x^2=-16$: $\res{\text{no té solució}}$
  \task $x^{2}-4x=0$: $\res{x=0,\ \ x=4}$
  \task $2x^{2}-7x=0$: $\res{x=0,\ \ x=\frac{7}{2}}$
  \task $x^2=\frac{25}{4}$: $\res{x=\pm\frac{5}{2}}$
  \task $x^2=5$: $\res{x=\pm\sqrt{5}}$
  \task $x^{2}-6x=0$: $\res{x=0,\ \ x=6}$
  \task $x^2=-\frac{4}{9}$: $\res{\text{no té solució}}$
\end{tasks}

\ex{2} \begin{tasks}(2)
  \task $x^2=9$: $\res{x=\pm3}$
  \task $x^{2}+6x=0$: $\res{x=-6,\ \ x=0}$
  \task $x^2=4$: $\res{x=\pm2}$
  \task $x^2=9$: $\res{x=\pm3}$
  \task $2x^{2}-3x=0$: $\res{x=0,\ \ x=\frac{3}{2}}$
  \task $x^2=40$: $\res{x=\pm2 \sqrt{10}}$
\end{tasks}

\ex{3} $x^2=-c$ té solucions si $\res{c\le0}$ (si $c=0$, només $x=0$). \ $x(2x+b)=0$: dues solucions si $b\neq0$ ($x=0$ i $x=-\frac b2$) i una (doble) si $b=0$.


% =====================================================
\seccio{5}{Equacions de segon grau completes}

\exf{1} \begin{tasks}(3)
  \task $\res{x=3,\ \ x=5}$
  \task $\res{x=-5,\ \ x=2}$
  \task $\res{x=\frac{1}{2},\ \ x=2}$
  \task $\res{x=5\ \text{(doble)}}$
  \task $\res{\text{no té solució}}$
  \task $\res{x=-\frac{2}{3},\ \ x=\frac{1}{2}}$
\end{tasks}

\ex{2} \begin{tasks}(3)
  \task $\res{x=2\pm \sqrt{3}}$
  \task $\res{x=-1\pm \sqrt{5}}$
  \task $\res{x=\frac{3\pm \sqrt{3}}{2}}$
\end{tasks}

\ex{2} \begin{tasks}(2)
  \task $x^{2}-3x+2=0$: $\res{x=1,\ \ x=2}$
  \task $x^{2}-4x+4=0$: $\res{x=2\ \text{(doble)}}$
  \task $x^{2}+3x-10=0$: $\res{x=-5,\ \ x=2}$
\end{tasks}

\ex{2} \begin{tasks}(3)
  \task $\Delta=-31$: $\res{\text{cap}}$
  \task $\Delta=0$: $\res{\text{una (doble)}}$
  \task $\Delta=13$: $\res{\text{dues}}$
\end{tasks}

\ex{3} $\Delta=k^2-36=0\Rightarrow\res{k=\pm6}$; \ amb $k=6$, $\res{x=-3}$; amb $k=-6$, $\res{x=3}$.

\ex{3} $(x-2)(2x+3)=0\Rightarrow\res{2x^2-x-6=0}$ (o qualsevol múltiple).


% =====================================================
\seccio{6}{Segon grau amb parèntesis i denominadors}

\exf{2} \begin{tasks}(2)
  \task $x^{2}-8x+15=0$: $\res{x=3,\ \ x=5}$
  \task $x^{2}-3x+2=0$: $\res{x=1,\ \ x=2}$
  \task $x^{2}-x-6=0$: $\res{x=-2,\ \ x=3}$
  \task $x^{2}-2x-8=0$: $\res{x=-2,\ \ x=4}$
\end{tasks}

\ex{2} \begin{tasks}(1)
  \task $x^{2}-x-2=0$: $\res{x=-1,\ \ x=2}$
  \task $2x^{2}-x-6=0$: $\res{x=-\frac{3}{2},\ \ x=2}$
  \task $x^{2}-4x+3=0$: $\res{x=1,\ \ x=3}$
  \task $x^{2}-6x+8=0$: $\res{x=2,\ \ x=4}$
\end{tasks}

\ex{3} \begin{tasks}(1)
  \task $x^{2}+3x-10=0$: $\res{x=-5,\ \ x=2}$
  \task $5x^{2}-9x-2=0$: $\res{x=-\frac{1}{5},\ \ x=2}$
  \task $x^{2}-3x=0$: $\res{x=0,\ \ x=3}$
  \task $x^{2}-2x-6=0$: $\res{x=1\pm \sqrt{7}}$
\end{tasks}


% =====================================================
\seccio{7}{Problemes de segon grau}

\exf{2} $x^2+x-110=0\Rightarrow x=10,\ x=-11$; ha de ser natural: $\res{x=10}$.

\ex{2} $x(x+2)=168\Rightarrow x^2+2x-168=0\Rightarrow x=12,\ x=-14$: $\res{12\ \text{i}\ 14\ \text{anys}}$

\ex{2} Els costats sumen 17: $x(17-x)=60\Rightarrow x^2-17x+60=0\Rightarrow x=12,\ x=5$: $\res{12\ \text{m}\times5\ \text{m}}$

\ex{2} $x^2+(x+7)^2=169\Rightarrow x^2+7x-60=0\Rightarrow x=5,\ x=-12$: $\res{5\ \text{cm}\ \text{i}\ 12\ \text{cm}}$; $\res{A=60\ \text{cm}^2}$

\ex{2} $(x+3)(x-2)=66\Rightarrow x^2+x-72=0\Rightarrow x=8,\ x=-9$: $\res{8\ \text{m}}$

\ex{3} $(20+2x)(15+2x)=594\Rightarrow4x^2+70x-294=0\Rightarrow2x^2+35x-147=0\Rightarrow x=\dfrac{-35\pm49}{4}$: $\res{x=3,5\ \text{cm}}$ (la negativa no val).

\ex{3} $\dfrac{120}{x}-\dfrac{120}{x+2}=5\Rightarrow240=5x(x+2)\Rightarrow x^2+2x-48=0\Rightarrow x=6,\ x=-8$: $\res{6\ \text{amics}}$ (20 € cadascú).

\ex{3} $x^2+(x+1)^2=(x+2)^2\Rightarrow x^2-2x-3=0\Rightarrow x=3,\ x=-1$: $\res{3,\ 4\ \text{i}\ 5}$


% =====================================================
\seccio{8}{Equacions factoritzades}

\exf{1} \begin{tasks}(3)
  \task $\res{x=-5,\ \ x=3}$
  \task $\res{x=0,\ \ x=\frac{7}{2}}$
  \task $\res{x=-\frac{4}{3},\ \ x=-1,\ \ x=2}$
  \task $\res{x=6}$
  \task $\res{x=-4,\ \ x=0,\ \ x=\frac{1}{4}}$
  \task $\res{x=-3,\ \ x=1,\ \ x=3}$
\end{tasks}

\ex{2} \begin{tasks}(2)
  \task $\res{x=-4,\ \ x=2,\ \ x=3}$
  \task $\res{x=-4,\ \ x=4}$ ($x^2+4=0$ no en dona)
  \task $\res{x=-2,\ \ x=0,\ \ x=5}$
  \task $\res{x=- \sqrt{2},\ \ x=\frac{1}{2},\ \ x=\sqrt{2}}$
\end{tasks}

\ex{2} \begin{tasks}(2)
  \task $x(x+3)(x-3)=0$: $\res{x=-3,\ \ x=0,\ \ x=3}$
  \task $2x^2(x-4)=0$: $\res{x=0,\ \ x=4}$
  \task $x(x-5)^2=0$: $\res{x=0,\ \ x=5}$
  \task $x^2(x+1)(x-1)=0$: $\res{x=-1,\ \ x=0,\ \ x=1}$
  \task $x(x+3)(x-2)=0$: $\res{x=-3,\ \ x=0,\ \ x=2}$
  \task $3x(x+2)(x-2)=0$: $\res{x=-2,\ \ x=0,\ \ x=2}$
\end{tasks}

\ex{3} \begin{tasks}(1)
  \task $(x-1)(x+2)-(x-1)=0\Rightarrow(x-1)(x+1)=0$: $\res{x=-1,\ \ x=1}$
  \task $x(x-3)-2(x-3)=0\Rightarrow(x-3)(x-2)=0$: $\res{x=2,\ \ x=3}$
\end{tasks}


% =====================================================
\seccio{9}{Equacions racionals}

\exf{1} \begin{tasks}(2)
  \task m.c.m.$=x$: $x-1=0$: $\res{x=1}$
  \task m.c.m.$=x - 2$: $x-3=0$: $\res{x=3}$
  \task m.c.m.$=x - 1$: $x-2=0$: $\res{x=2}$
  \task m.c.m.$=(x - 1) (x + 1)$: $x+5=0$: $\res{x=-5}$
  \task m.c.m.$=4 x$: $x^{2}-36=0$: $\res{x=-6,\ \ x=6}$
  \task m.c.m.$=x (x + 2)$: $x-1=0$: $\res{x=1}$
\end{tasks}

\ex{2} \begin{tasks}(1)
  \task m.c.m.$=x$: $x^{2}-6x+8=0$: $\res{x=2,\ \ x=4}$
  \task m.c.m.$=2 x (x + 3)$: $x^{2}-x-6=0$: $\res{x=-2,\ \ x=3}$
  \task m.c.m.$=x (x - 1)$: $2x^{2}-5x+2=0$: $\res{x=\frac{1}{2},\ \ x=2}$
  \task m.c.m.$=(x - 1) (x + 1)$: $2x^{2}-5x-3=0$: $\res{x=-\frac{1}{2},\ \ x=3}$
\end{tasks}

\ex{3} \begin{tasks}(1)
  \task m.c.m.$=(x - 1) (x + 1)$: $x^{2}-x-2=0$; $x=-1$ anul·la un denominador: $\res{x=2}$
  \task m.c.m.$=x - 3$: $x-3=0$; $x=3$ anul·la un denominador: $\res{\text{no té solució}}$
  \task m.c.m.$=(x - 2) (x + 2)$: $x-3=0$: $\res{x=3}$
  \task m.c.m.$=6 x (x + 1)$: $5x^{2}-7x-6=0$: $\res{x=-\frac{3}{5},\ \ x=2}$
\end{tasks}

\ex{3} $x+\dfrac1x=\dfrac{10}{3}$; m.c.m.$=3 x$: $3x^{2}-10x+3=0$: $\res{x=\frac{1}{3},\ \ x=3}$ (dues possibilitats).


% =====================================================
\seccio{10}{Equacions amb radicals}

\exf{1} \begin{tasks}(3)
  \task $x+3=16$: $\res{x=13}$
  \task $2x-1=9$: $\res{x=5}$
  \task $x^2=16$: $\res{x=\pm4}$
  \task Una arrel no pot ser negativa: $\res{\text{no té solució}}$
  \task $\sqrt x=3$: $\res{x=9}$
  \task $5x-1=4$: $\res{x=1}$
\end{tasks}

\ex{2} \begin{tasks}(1)
  \task $x+7=x^2+2x+1\Rightarrow x^2+x-6=0\Rightarrow x=2,\ x=-3$; $x=-3$ no val ($\sqrt4\neq-2$): $\res{x=2}$
  \task $x-4=\sqrt{x+2}\Rightarrow x^2-9x+14=0\Rightarrow x=7,\ x=2$; $x=2$ no val ($2-2\neq4$): $\res{x=7}$
  \task $\sqrt{2x+1}=x-1\Rightarrow x^2-4x=0\Rightarrow x=0,\ x=4$; $x=0$ no val ($1+1\neq0$): $\res{x=4}$
  \task $3x-2=x^2\Rightarrow x^2-3x+2=0$; totes dues valen: $\res{x=1,\ \ x=2}$
  \task $\sqrt{x-1}=3-x\Rightarrow x^2-7x+10=0\Rightarrow x=2,\ x=5$; $x=5$ no val ($5+2\neq3$): $\res{x=2}$
  \task $4(x+1)=(x+1)^2\Rightarrow x^2-2x-3=0$; totes dues valen: $\res{x=-1,\ \ x=3}$
\end{tasks}

\ex{2} \begin{tasks}(1)
  \task $2x+3=x+5$: $\res{x=2}$
  \task $5x-4=x+8$: $\res{x=3}$
  \task $9(x-1)=x+7$: $\res{x=2}$
  \task $x^2-x-12=0$; totes dues valen ($\sqrt9=\sqrt9$, $\sqrt2=\sqrt2$): $\res{x=4,\ \ x=-3}$
  \task $x^2+3x=4x\Rightarrow x^2-x=0$; totes dues valen: $\res{x=0,\ \ x=1}$
  \task $x^2-2x-3=0\Rightarrow x=3,\ x=-1$; $x=-1$ no val ($\sqrt{-2}$ no existeix): $\res{x=3}$
\end{tasks}

\ex{3} \begin{tasks}(1)
  \task $\sqrt{25-x^2}=x-1\Rightarrow 25-x^2=x^2-2x+1\Rightarrow x^2-x-12=0\Rightarrow x=4,\ x=-3$; $x=-3$ no val ($-3-4\neq1$): $\res{x=4}$
  \task $\sqrt{x^2+5}=2x-1\Rightarrow x^2+5=4x^2-4x+1\Rightarrow3x^2-4x-4=0\Rightarrow x=2,\ x=-\frac23$; $x=-\frac23$ no val ($\frac73+1\neq-\frac43$): $\res{x=2}$
\end{tasks}

\ex{3} $\sqrt{x+2}=x-4\Rightarrow x^2-9x+14=0\Rightarrow x=7,\ x=2$; $x=2$ no val: $\res{x=7}$


\end{document}
